\section{Rangkuman}
\label{sec:rangkuman-bab4}

Pada bab ini telah dibahas:

\begin{itemize}
    \item Mikroprosesor berkembang dari Intel 4004 (4 bit) hingga
          8086 (16 bit) yang menjadi dasar keluarga x86. Varian 8088
          dengan bus luar 8 bit dipakai pada IBM PC.

    \item Arsitektur dipandang dari dua sisi: \term{internal}
          (register, ALU, dan unit di dalam chip) dan
          \term{eksternal} (pin, bus, dan sinyal kendali).

    \item 8086 terbagi menjadi \eng{Bus Interface Unit} (BIU) dan
          \eng{Execution Unit} (EU). BIU memegang seluruh lalu lintas
          bus luar; EU hanya menjalankan instruksi dan tidak
          terhubung langsung ke pin luar.

    \item BIU memuat register segmen (\reg{CS}, \reg{DS}, \reg{SS},
          \reg{ES}), \reg{IP}, penghitung alamat 20-bit, dan
          \eng{instruction queue} enam byte.

    \item EU memuat \eng{ALU}, register umum, register penunjuk dan
          indeks, serta register \eng{flag}.

    \item Rumus alamat fisik adalah (segmen $\times$ 16) $+$ offset,
          sehingga 8086 dapat menjangkau 1 MB memori meskipun
          register-nya hanya 16 bit.

    \item \eng{Instruction queue} memungkinkan \eng{pipelining}:
          pengambilan instruksi berikutnya berimpit dengan eksekusi
          instruksi sekarang. Instruksi lompatan mengosongkan
          antrean (\eng{queue flush}) dan menghapus keuntungan ini.

    \item Tiga kelompok bus eksternal adalah bus alamat 20-bit
          (satu arah), bus data 16-bit (dua arah), dan bus kontrol.
          Pin \reg{AD0}--\reg{AD15} dimultipleks sehingga alamat
          harus ditangkap dengan bantuan sinyal \reg{ALE}.

    \item Siklus instruksi terdiri atas \eng{fetch}, \eng{decode},
          \eng{execute}, dan \eng{store}. Tahap \eng{store} hanya
          memerlukan siklus bus bila tujuannya berada di memori.
\end{itemize}

\begin{checklistbox}
Tandai setiap butir yang sudah Anda kuasai.

\begin{itemize}
    \item[\checkbox] Saya dapat menjelaskan perbedaan arsitektur
          internal dan eksternal sebuah mikroprosesor.
    \item[\checkbox] Saya dapat menggambar diagram blok BIU/EU
          beserta isi register masing-masing unit.
    \item[\checkbox] Saya dapat menjelaskan peran
          \eng{instruction queue} enam byte dan \eng{pipelining}.
    \item[\checkbox] Saya dapat menghitung alamat fisik dari pasangan
          segmen dan offset.
    \item[\checkbox] Saya dapat menyebutkan fungsi bus alamat, bus
          data, bus kontrol, dan sinyal \reg{ALE}, \code{M/IO},
          \reg{RD}, \reg{WR}, \reg{READY}, \reg{RESET}.
    \item[\checkbox] Saya dapat menelusuri satu instruksi melalui
          tahapan \eng{fetch}--\eng{decode}--\eng{execute}--\eng{store}.
    \item[\checkbox] Saya dapat menjelaskan mengapa instruksi yang
          mengakses memori lebih lambat daripada instruksi yang
          hanya memakai register.
\end{itemize}

\vspace{0.4em}
Bila masih ada butir yang belum bertanda, ulangi bagian yang
bersesuaian sebelum melanjutkan ke bab berikutnya.
\end{checklistbox}

\begin{infobox}
\textbf{Rujukan bab.} Materi pada bab ini diperkuat oleh
\cite{intel1979} sebagai rujukan utama datasheet dan pembagian
BIU/EU, \cite{stallings2016} untuk gambaran organisasi dan
arsitektur komputer, \cite{courses2103} untuk penjelasan segmentasi
dan alamat fisik, serta \cite{carleton2011} untuk konteks
keluarga x86. Untuk penelusuran lebih dalam mengenai bus dan sinyal
kendali 8086, lihat \cite{intel1979}.
\end{infobox}
